\section*{Journal Section3: Structure theorem for quadratics}
\textit{Shir Peleg and Amir Shpilka, Forum of Mathematics, Sigma10:e112 (2022), pp11--16, DOI10.1017/fms.2022.100. CC BY4.0. The numbering below follows the journal version.}
\noindent\textit{Editorial reference-number concordance: journal reference [12] is Gupta (2014), listed as CCC reference [17] in this material. Journal Theorem2.11 is the resultant criterion stated as CCC Theorem17 in the preceding context. The journal continuation retains its original numbering.}
\begin{nttheorem}[Author Theorem3.1]
Let $\{Q_k\}_{k\in\mathcal K},A,B$ be homogeneous polynomials of degree $2$ such that $\prod_{k\in\mathcal K}Q_k\in\sqrt{\langle A,B\rangle}$. Then one of the following cases hold:
\begin{description}
\item[(prime-case)] There is $k\in\mathcal K$ such that $Q_k$ is in the linear span of $A,B$.
\item[(product-case)] There exists a nontrivial linear combination of the form $\alpha A+\beta B=c\cdot d$, where $c$ and $d$ are linear forms.
\item[(linear-case)] There exist two linear forms $c$ and $d$ such that when setting $c=d=0$ we get that $A,B$ and one of $\{Q_k\}_{k\in\mathcal K}$ vanish.
\end{description}
\end{nttheorem}
The following claim of [12] shows that we can assume $|\mathcal K|=4$ in the statement of Theorem3.1.
\begin{ntclaim}[Author Claim3.2, Claim11 in journal reference12]
Let $P_1,\ldots,P_d,Q_1,\ldots,Q_k\in\mathbb C[x_1,\ldots,x_n]$ be homogeneous and the degree of each $P_i$ is at most $r$. Then,
\[\prod_{i=1}^kQ_i\in\sqrt{\langle P_1,\ldots,P_d\rangle}\ \Rightarrow\ \exists\{i_1,\ldots,i_{r^d}\}\subset[k]\text{ such that }\prod_{j=1}^{r^d}Q_{i_j}\in\sqrt{\langle P_1,\ldots,P_d\rangle}.\]
\end{ntclaim}
Thus, for $r=d=2$ it follow that there are at most four polynomials among the $Q_i$s whose product is in $\sqrt{\langle A,B\rangle}$.
Before proving Theorem3.1, we explain the intuition behind the different cases in the theorem. Clearly, if one of $Q_1,\ldots,Q_4$ is a linear combination of $A,B$, then it is in their radical (and in fact, in their linear span). If $A$ and $B$ span a product of the form $ab$, then, say, $(A+ac)(A+bd)$ is in their radical. Indeed, $\sqrt{\langle A,B\rangle}=\sqrt{\langle A,ab\rangle}$. This case is clearly different than the linear span case. Finally, we note that, if $A=ac+bd$ and $B=ae+bf$, then the product $a\cdot b\cdot(cf-de)$ is in $\sqrt{\langle A,B\rangle}$.\footnote{If we insist on having all factors of degree2, then the same argument shows that the product $(a^2+A)\cdot(b^2+B)\cdot(cf-de)$ is in $\sqrt{\langle A,B\rangle}$.} This case is different than the other two cases as $A$ and $B$ do not span any linear form (or any reducible quadratic) non trivially.
Thus, all the three cases are distinct and can happen. What Theorem3.1 shows is that, essentially, these are the only possible cases.
\begin{proof}[Author proof of Theorem3.1]
Following Theorem3.2, we shall assume in the proof that $|\mathcal K|=4$. By applying a suitable linear transformation, we can assume that for some $r\geq1$
\[A=\sum_{i=1}^rx_i^2.\]
We can also assume without loss of generality that $x_1^2$ appears only in $A$ as we can replace $B$ with any polynomial of the form $B'=B-\alpha A$ without affecting the result as $\langle A,B\rangle=\langle A,B'\rangle$. Furthermore, all cases in the theorem remain the same if we replace $B$ with $B'$ and vice versa.
In a similar fashion, we can replace $Q_1$ with $Q'_1=Q_1-\alpha A$ to get rid of the term $x_1^2$ in $Q_1$. We can do the same for the other $Q_i$s. Thus, without loss of generality, the situation is
\begin{align}A&=x_1^2-A'\notag\\B&=x_1\cdot b-B'\tag{3.1}\\Q_i&=x_1\cdot b_i-Q'_i\quad\text{for }i\in\{1,2,3,4\},\notag\end{align}
where $A',b,B',Q'_i,b_i$ are homogeneous polynomials that do not depend on $x_1$, the polynomials $A',B',Q'$ have degree $\leq2$ and $b,b_i$ are linear forms. The analysis shall deal with two cases according to whether $B$ depends on $x_1$ or not, as we only consider the resultant of $A$ and $B$ with respect to $x_1$ when it appears in both polynomials.
\paragraph{Case $b\not\equiv0$}
Consider the resultant of $A$ and $B$ with respect to $x_1$. It is easy to see that
\[\operatorname{Res}_{x_1}(A,B)=B'^2-b^2\cdot A'.\]
We first prove that, if the resultant is irreducible, then Case(prime-case) of Theorem3.1 holds. For this, we shall need the following claim.
\begin{ntclaim}[Author Claim3.3]
Whenever $\operatorname{Res}_{x_1}(A,B)=0$, it holds that $\prod_{i=1}^4(B'\cdot b_i-b\cdot Q'_i)=0$.
\end{ntclaim}
\begin{proof}
Let $\boldsymbol\alpha\in\mathbb C^{n-1}$ be such that $\operatorname{Res}_{x_1}(A,B)(\boldsymbol\alpha)=0$ then either $b(\boldsymbol\alpha)=0$, which also implies $B'(\boldsymbol\alpha)=0$ and in this case the claim clearly holds, or $b(\boldsymbol\alpha)\ne0$. Consider the case $b(\boldsymbol\alpha)\ne0$, and set $x_1=B'(\boldsymbol\alpha)/b(\boldsymbol\alpha)$ (we are free to select a value for $x_1$ as $\operatorname{Res}_{x_1}(A,B)$ does not involve $x_1$). Notice that for this substitution we have that $B(\boldsymbol\alpha)=0$ and that
\[A|_{x_1=B'(\boldsymbol\alpha)/b(\boldsymbol\alpha)}=(B'(\boldsymbol\alpha)/b(\boldsymbol\alpha))^2-A'(\boldsymbol\alpha)=\operatorname{Res}_{x_1}(A,B)(\boldsymbol\alpha)/b(\boldsymbol\alpha)^2=0.\]
Hence, we also have $\prod_{i=1}^4Q_i|_{x_1=B'(\boldsymbol\alpha)/b(\boldsymbol\alpha)}=0$. In other words that
\[\left(\frac1{b^4}\prod_{i=1}^4(B'\cdot b_i-b\cdot Q'_i)\right)(\boldsymbol\alpha)=0.\]
\end{proof}
It follows that
\[\prod_{i=1}^4(B'\cdot b_i-b\cdot Q'_i)\in\sqrt{\langle\operatorname{Res}_{x_1}(A,B)\rangle}.\]
In other words, for some positive integer $k$ we have that $\operatorname{Res}_{x_1}(A,B)$ divides $\left(\prod_{i=1}^4(B'\cdot b_i-b\cdot Q'_i)\right)^k$. As every irreducible factor of $\left(\prod_{i=1}^4(B'\cdot b_i-b\cdot Q'_i)\right)^k$ is of degree3 or less, we get that, if the resultant is irreducible, then one of the multiplicands must be identically zero. Assume without loss of generality that $B'b_1-bB'_1=0$. It is not hard to verify that in this case either $Q_1$ is a scalar multiple of $B$ and then Theorem3.1(prime-case) holds, or that $B'$ is divisible by $b$. However, in the latter case it also holds that $b$ divides the resultant, contradicting the assumption that it is irreducible.
From now on we assume that $\operatorname{Res}_{x_1}(A,B)$ is reducible. We consider two possibilities. Either $\operatorname{Res}_{x_1}(A,B)$ has a linear factor or it can be written as
\[\operatorname{Res}_{x_1}(Q_1,Q_2)=C\cdot D,\]
for irreducible quadratic polynomials $C$ and $D$.
Consider the case where the resultant has a linear factor. If that linear factor is $b$, then $b$ also divides $B$ and Theorem3.1(product-case) holds. Otherwise, if it is a different linear form $\ell$, then when setting $\ell=0$ we get that the resultant of $A|_{\ell=0}$ and $B|_{\ell=0}$ is zero, and hence, either $B|_{\ell=0}$ is identically zero and Theorem3.1(product-case) holds or they share a common factor (see Theorem2.11). It is not hard to see that, if that common factor is of degree2, then Theorem3.1(product-case) holds, and if it is a linear factor, then Theorem3.1(linear-case) holds.
Thus, the only case left to handle (when $b\not\equiv0$) is when there are two irreducible quadratic polynomials, $C$ and $D$, such that $CD=\operatorname{Res}_{x_1}(A,B)$. As $C$ and $D$ divide two multiplicands in $\prod_{i=1}^4(B'\cdot b_i-b\cdot Q'_i)$, we can assume, without loss of generality, that $(B'\cdot b_3-b\cdot Q'_3)\cdot(B'\cdot b_4-b\cdot Q'_4)\in\sqrt{\langle\operatorname{Res}_{x_1}(A,B)\rangle}$. Next, we express $A',B',C$ and $D$ as quadratics over $b$. That is,
\begin{align}A'&=\alpha b^2+a_1b+A''\tag{3.2}\\B'&=\beta b^2+a_2b+B''\notag\\C&=\gamma b^2+a_3b+C''\notag\\D&=\delta b^2+a_4b+D'',\notag\end{align}
where $a_1,\ldots,D''$ do not involve $b$ (nor $x_1$). We have the following two representations of the resultant:
\begin{align}\operatorname{Res}_{x_1}(A,B)&=B'^2-b^2\cdot A'\tag{3.3}\\&=\beta^2b^4+2\beta a_2b^3+(2\beta B''+a_2^2)b^2+2a_2B''b+B''^2-\alpha b^4-a_1b^3-A''b^2\notag\\&=(\beta^2-\alpha)b^4+(2\beta a_2-a_1)b^3+(2\beta B''+a_2^2-A'')b^2+2a_2B''b+B''^2\notag\end{align}
and
\begin{align}\operatorname{Res}_{x_1}(Q_1,Q_2)&=CD\tag{3.4}\\&=(\gamma b^2+a_3b+C'')(\delta b^2+a_4b+D'')\notag\\&=\gamma\delta b^4+(\gamma a_4+\delta a_3)b^3+(\gamma D''+a_3a_4+\delta C'')b^2+(a_3D''+a_4C'')b+C''D''.\notag\end{align}
Comparing the different coefficients of $b$ in the two representations in equations(3.3) and(3.4), we obtain the following equalities:
\begin{equation}B''^2=C''D''\tag{3.5}\end{equation}
\begin{equation}2a_2B''=a_3D''+a_4C''.\tag{3.6}\end{equation}
